\subsection{Analisis Kanal DRUMS dan VOCALS dalam Format WAV}
Analisis tambahan ini membandingkan dua kanal audio yang sudah
tersedia, khususnya pola kejadian singkat dan komponen yang lebih
jarang muncul. Sumber kedua kanal berbeda dari lima instrumen
SoundFont, sehingga perolehan data dijelaskan sebelum hasilnya.

\subsubsection{Konversi Format dan Validasi Sampel}
Channel 11 dan 12 pada FLP masing-masing merujuk ke
\texttt{DRUMS.mp3} dan \texttt{VOCALS.mp3}. Tidak terdapat not pattern
untuk merender ulang kedua kanal dengan MIDI. Karena itu, audio
didekode lalu disimpan sebagai WAV float32 tanpa resampling,
normalisasi, atau perubahan amplitudo menggunakan SoundFile
\cite{soundfile}. Masing-masing berisi 7.090.031 frame stereo pada
44.100 Hz, berdurasi 160,771678 s.

MP3 merupakan sumber yang telah mengalami kompresi lossy. WAV hasil
konversi hanya menyimpan hasil dekodenya tanpa kompresi lossy baru;
detail yang hilang tidak dipulihkan. Pemeriksaan baca-tulis
menunjukkan kesamaan setiap sampel float32, dengan galat maksimum
nol. Kedua WAV tidak senyap dan tidak memiliki sampel dengan
magnitudo mencapai satu. Puncak amplitudo stereo DRUMS sekitar
0,89477, sedangkan VOCALS 0,99853.

Konsekuensi matematisnya sederhana: jika vektor sampel sama,
RMS, DFT dan STFT dengan parameter sama juga sama. Analisis DRUMS
dihitung ulang dari WAV; RMS, energi dan puncak frekuensinya sama
dengan perhitungan langsung dari dekode MP3. Dengan demikian, penggunaan WAV
di sini adalah pilihan format kerja, bukan peningkatan ketelitian
analisis atau kualitas sumber.

Parameter stretching yang terbaca pada kedua channel adalah
multiplier satu, pitch nol dan time nol, dengan start offset
parameter nol. Namun, pemetaan offset klip audio ke sampel serta
hasil playback belum diperiksa terhadap render FL Studio acuan.
Analisis memakai berkas penuh, tanpa memotong atau memaksanya menjadi
69,23 s. Kedua kanal tidak dimasukkan ke
\texttt{All\_Instruments.wav}. Nama VOCALS mengikuti label sumber
dan tidak dianggap sebagai bukti bahwa semua komponennya suara
manusia murni.

\subsubsection{Aktivitas Temporal dan Amplitudo Efektif}
Perbandingan waktu bertujuan mengetahui apakah kedua kanal aktif
sepanjang rekaman atau hanya pada beberapa bagian. Waveform dan
RMS dihitung dengan metode yang sama seperti instrumen rekonstruksi.
\begin{figure}[H]
\centering\includegraphics[width=\linewidth]{01_waveform_rms_kanal.png}
\caption{Amplitudo dan RMS kanal audio penuh. X: waktu dari awal
berkas sumber; kiri: amplitudo digital; kanan: RMS 100 ms.
Skala antarbaris sama.}\label{fig:kanalwaktu}
\end{figure}
Gambar~\ref{fig:kanalwaktu} memperlihatkan DRUMS aktif pada sebagian
besar rekaman, sedangkan VOCALS mempunyai kelompok kejadian utama
sekitar 31,4--35,2 s, 92,3--100,0 s dan 153,0--155,9 s.
Waktu ini mengikuti batas blok aktivitas, bukan transkripsi suku
kata atau waktu onset yang persis. Dengan ambang 1\% RMS blok
maksimum tiap kanal, durasi aktif DRUMS sekitar 156,1 s dan VOCALS
14,2 s.

RMS global VOCALS sekitar 0,04895, tetapi RMS pada blok aktif sekitar
0,16470. Untuk DRUMS, keduanya lebih dekat, yaitu 0,07757 dan 0,07872.
Bagian diam menurunkan rerata kuadrat sepanjang rekaman. Karena itu,
RMS global VOCALS yang lebih kecil tidak berarti setiap kejadiannya
beramplitudo lebih kecil daripada DRUMS. Perubahan dari latar
beramplitudo rendah menuju kejadian kuat juga mendukung deskripsi
nonstasioner, tanpa menyatakan hasil uji statistik formal.

\subsubsection{Spektrum Global dan Distribusi Energi}
FFT digunakan untuk membandingkan wilayah frekuensi yang paling
berkontribusi pada kedua kanal. Semua sampel berkas penuh digunakan;
jarak bin sekitar 0,006220 Hz. Grafik berikut memperlihatkan
0--5000 Hz, sedangkan perhitungan energi tetap sampai Nyquist.
\begin{figure}[H]
\centering\includegraphics[width=.96\linewidth]{02_fft_kanal.png}
\caption{Kerapatan daya DRUMS dan VOCALS dengan acuan sama.
X: frekuensi; Y: daya per Hz dalam dB. Tampilan merangkum maksimum
per 0,5 Hz; perhitungan memakai bin asli.}\label{fig:kanalfft}
\end{figure}
\begin{table}[H]
\centering\small
\caption{Hasil perhitungan dari dua WAV hasil dekode.}
\begin{tabular}{lrrrrr}
\toprule
Kanal & RMS & Puncak (Hz) & Sentroid (Hz) & $<250$ Hz (\%) & $\geq1000$ Hz (\%) \\
\midrule
DRUMS & 0.07757 & 41.60 & 563.50 & 91.19 & 6.83 \\
VOCALS & 0.04895 & 780.08 & 1195.31 & 0.02 & 31.87 \\
\bottomrule
\end{tabular}

\end{table}
DRUMS memiliki sekitar 91,19\% energi di bawah 250 Hz, sedangkan
VOCALS hanya sekitar 0,017\%. Pada VOCALS, sekitar 68,11\% berada
pada 250--1000 Hz dan 31,87\% pada 1000 Hz ke atas. Komposisi
VOCALS lebih banyak berkontribusi di wilayah menengah dan tinggi
dibandingkan DRUMS. Ini perbedaan distribusi spektral, bukan bukti
perbedaan keras suara perseptual atau kemurnian jenis sumber.

Puncak terpilih VOCALS sekitar 780,08 Hz tidak otomatis merupakan
frekuensi dasar suara manusia. Komponen yang bertahan, harmonik,
dan perubahan nada sama-sama berkontribusi pada FFT global.
Sentroid VOCALS sekitar 1195,31 Hz, dibandingkan DRUMS 563,50 Hz,
menunjukkan pusat distribusi daya pada frekuensi lebih tinggi.
Parseval kedua WAV menghasilkan galat relatif sekitar
$2,1\times10^{-15}$ atau lebih kecil. Hasil DRUMS lama tetap
tersimpan dan tidak digunakan sebagai pengganti perhitungan VOCALS.

\clearpage
\subsubsection{Karakteristik Waktu--Frekuensi DRUMS}
STFT dipakai untuk melokalisasi kejadian pada spektrum DRUMS.
Gambar~\ref{fig:stftdrumswav} dan~\ref{fig:stftvocalswav} menggunakan
29--33 s menurut waktu berkas sumber, frekuensi 0--5000 Hz, serta
skala warna yang sama, $-110$ sampai $-20$ dB. Kedua panjang Hann
tetap 1024 dan 4096, dengan overlap 75\%. Batas tampilan ini berbeda
dari panel lima instrumen; perbandingan warna langsung dilakukan
di dalam pasangan yang memakai skala sama.
\begin{figure}[H]
\centering\includegraphics[width=\linewidth]{stft_zoom_DRUMS.png}
\caption{STFT DRUMS dari WAV hasil dekode, pada waktu sumber 29--33 s.
Atas: 1024 sampel; bawah: 4096 sampel.}\label{fig:stftdrumswav}
\end{figure}
X adalah pusat jendela; Y adalah frekuensi. Daerah terang rendah
disertai guratan yang menjangkau frekuensi lebih tinggi, terutama
pada kejadian berulang sekitar 29--31 s. Pola itu sesuai dengan
perubahan amplitudo singkat yang melibatkan banyak frekuensi,
bukan hanya satu osilasi tetap. Panel 4096 memperinci garis
horizontal, sementara panel 1024 melibatkan potongan waktu lebih
pendek. Perkalian dengan Hann menjelaskan kompromi antara
penyebaran frekuensi dan dukungan waktu. Grafik ini tidak dipakai
untuk mengidentifikasi jenis setiap pukulan drum secara otomatis.

\clearpage
\subsubsection{Karakteristik Waktu--Frekuensi VOCALS}
STFT VOCALS memperlihatkan hubungan antara periode beramplitudo
rendah dan kejadian yang membentuk puncak pada FFT global.
\begin{figure}[H]
\centering\includegraphics[width=\linewidth]{stft_zoom_VOCALS.png}
\caption{STFT VOCALS pada waktu sumber 29--33 s. Komponen kuat mulai
tampak sekitar 31,4 s. Nama kanal mengikuti sumber FLP.}
\label{fig:stftvocalswav}
\end{figure}
Sebelum sekitar 31,4 s, sebagian besar panel gelap pada skala yang
dipilih. Setelah itu muncul jalur terang sekitar 0,8 kHz dan
kelipatannya, dengan perubahan posisi sebelum dan sesudah bagian
yang relatif mendatar sekitar 31,8--32,5 s. Hubungan mendekati
kelipatan tersebut mendukung keberadaan struktur harmonik, tetapi
tidak cukup untuk menetapkan fundamental atau mekanisme sumbernya.

Panel 4096 menampilkan garis lebih tipis; pada panel 1024 komponen
berdekatan lebih melebar. Pergeseran jalur memperlihatkan bahwa
komponen kuat tidak selalu berada pada frekuensi yang tetap.
Secara matematis, FFT global merangkum seluruh perubahan itu,
sedangkan STFT melokalisasinya melalui jendela. Kedua hasil berasal
dari sampel dekode yang sama, sehingga detail gambar tidak boleh
diartikan sebagai informasi audio baru yang muncul akibat konversi
MP3 ke WAV. Cuplikan kedua kanal tersedia dalam notebook untuk
pemeriksaan pendengaran oleh pembaca.
